\section{Gabarito comentado — Simulado 1}

\begin{longtable}{c c p{10.8cm}}
\toprule
Q & Gab. & Comentário \\
\midrule
\endfirsthead
\toprule Q & Gab. & Comentário \\
\midrule
\endhead
1 & B & ``Embora'' expressa concessão; ``apesar de'' preserva a relação e o resultado negativo. \\
2 & B & Oração relativa restritiva, sem vírgulas, delimita o subconjunto devolvido. \\
3 & C & ``Existir'' concorda com ``evidências''; haver/fazer impessoais ficam no singular. \\
4 & B & ``Porquanto'' pode assumir valor causal/explicativo próximo de ``porque''. \\
5 & C & ``Seu'' pode retomar equipe ou gestor sem contexto desambiguador. \\
6 & B & ``Referência a'' + artigo feminino ``a norma'' produz crase. \\
7 & C & Validar domínio por canal confiável e reportar phishing evita captura de credenciais. \\
8 & B & Necessidade = tratamento limitado ao mínimo necessário à finalidade. \\
9 & B & A LAI adota publicidade como regra e sigilo como exceção legal. \\
10 & B & Privacidade, segurança e finalidade continuam exigíveis em serviços digitais. \\
11 & B & Negação de universal é existencial: existe ao menos um que não foi aprovado. \\
12 & A & Contrapositiva de $P\rightarrow Q$ é $\neg Q\rightarrow\neg P$. \\
13 & C & $800\times1{,}15=920$ mil. \\
14 & C & Inclusão-exclusão: $120+80-40=160$. \\
15 & C & De $x-y=2$, $x=y+2$; substituindo: $3y+2=14$, $y=4$, $x=6$. \\
16 & B & Desconcentração distribui competência dentro da mesma pessoa jurídica. \\
17 & C & Autoexecutoriedade permite execução direta em hipóteses admitidas pelo ordenamento. \\
18 & B & Anulação decorre de ilegalidade; revogação, em regra, de mérito administrativo. \\
19 & B & A disciplina vigente da improbidade exige dolo e enquadramento legal; ilegalidade isolada não basta. \\
20 & B & Regra constitucional: responsabilidade objetiva do Estado perante a vítima, com análise de nexo/excludentes. \\
21 & B & Municípios integram a organização federativa brasileira e possuem autonomia constitucional. \\
22 & B & Direitos fundamentais não se restringem ao art. 5º. \\
23 & B & O Congresso exerce controle externo com auxílio do TCU. \\
24 & B & CPI não pratica atos submetidos à reserva jurisdicional. \\
25 & B & MP é função essencial à Justiça. \\
26 & B & Controle interno e externo têm posições e competências distintas, mas complementares. \\
27 & B & Auditoria operacional examina desempenho, economicidade, eficiência, eficácia e efetividade. \\
28 & B & Jurisdição una assegura acesso judicial contra lesão ou ameaça a direito. \\
29 & B & Parecer prévio subsidia o julgamento político competente das contas de governo. \\
30 & A & A IN nº 50/2017 disciplina tomada de contas especial no TCE-MA. \\
31 & A & A IN nº 82/2025 trata de emendas parlamentares, transparência e rastreabilidade. \\
32 & B & Regimento não pode contrariar a Lei Orgânica. \\
33 & B & Rastreabilidade é reconstrução da cadeia completa da emenda e sua execução. \\
34 & A & Guaxenduba pertence ao conflito contra os franceses. \\
35 & A & Balaiada: Período Regencial, forte participação popular e causas sociais/políticas. \\
36 & B & Parnaíba, Gurupi e Tocantins-Araguaia são destacados como bacias limítrofes. \\
37 & A & Itaqui e corredores ferroviários integram cadeias nacionais e internacionais de exportação. \\
38 & B & A LBI adota modelo relacional: impedimento em interação com barreiras. \\
39 & A & ODS 16 trata de paz, justiça e instituições eficazes, responsáveis e inclusivas. \\
40 & A & Ações afirmativas buscam corrigir desigualdades e promover igualdade material. \\
41 & D & Em /26 restam 6 bits: $2^6=64$ endereços no bloco. \\
42 & A & 802.1X + EAP + RADIUS é arquitetura típica de controle de acesso à rede. \\
43 & A & Estado preso ao nó quebra failover; externalização/replicação ou desenho stateless reduz o problema. \\
44 & A & RAID e snapshot local compartilham riscos e não cobrem exclusão, ransomware e desastre do storage. \\
45 & A & São as quatro condições clássicas de Coffman para deadlock. \\
46 & B & AD usa LDAP, Kerberos, DNS e outros componentes; não é sinônimo de LDAP. \\
47 & A & Dependência transitiva de atributo não chave viola 3FN. \\
48 & A & MVCC mantém versões para coordenar concorrência e reduzir bloqueios conforme implementação. \\
49 & A & Fatos registram medidas/eventos e referenciam dimensões. \\
50 & A & CDC captura mudanças para integração/streaming. \\
51 & A & Lineage documenta origem, transformações e destino. \\
52 & A & Sob partição, CAP explicita trade-off C versus A. \\
53 & A & Requisito não mensurável dificulta teste e aceite. \\
54 & A & Sprint Review inspeciona incremento e adapta backlog com stakeholders. \\
55 & A & Ciclo Red-Green-Refactor. \\
56 & A & Credencial admin permanente viola menor privilégio e segregação. \\
57 & A & Readiness indica aptidão para receber tráfego. \\
58 & A & SLO precisa de SLI/métrica e janela de avaliação. \\
59 & A & MFA fatigue explora repetição de prompts para induzir aprovação. \\
60 & A & PAM controla contas/acessos privilegiados. \\
61 & A & Assinatura digital apoia integridade, autenticidade e não repúdio; não cifra automaticamente conteúdo. \\
62 & A & Cópia imutável/offline reduz comprometimento simultâneo por ransomware. \\
63 & A & Cadeia de custódia documenta percurso e integridade da evidência. \\
64 & A & Zero Trust = verificar continuamente, contexto, identidade e menor privilégio. \\
65 & A & EDM é domínio de governança no COBIT. \\
66 & A & Governança avalia, direciona e monitora; gestão planeja/constrói/executa/monitora atividades. \\
67 & A & ITIL 4 enfatiza cocriar valor em relações de serviço. \\
68 & A & Probabilidade e impacto são componentes centrais da análise de risco. \\
69 & A & Gateways controlam ramificação e convergência de fluxos BPMN. \\
70 & A & KPI local pode incentivar comportamento desalinhado do resultado final. \\
71 & A & Essas são as três fases macro da IN SGD/ME 94/2022. \\
72 & A & Ignorar referências públicas comparáveis pode distorcer preço estimado e elevar risco de sobrepreço. \\
73 & A & Evidência produzida apenas pela parte interessada é pouco independente. \\
74 & A & Portabilidade, interoperabilidade e saída testada mitigam lock-in. \\
75 & B & 30 dias = 43.200 min; 0,1\% = 43,2 min. \\
76 & A & Dimensionar pelo total sem relação com demanda gera licenças ociosas e sobrecontratação. \\
77 & A & Em IaaS, cliente normalmente gerencia SO, middleware/aplicações e configurações sob sua responsabilidade. \\
78 & A & Elasticidade é ajuste dinâmico de capacidade. \\
79 & A & A divisão varia com IaaS/PaaS/SaaS e configuração concreta. \\
80 & A & Storage público por configuração é exemplo clássico de responsabilidade do consumidor. \\
81 & A & FinOps conecta custo, valor e decisões técnicas/financeiras. \\
82 & A & Multi-region aumenta complexidade de dados, consistência, custo e failover; exige teste. \\
83 & A & Mediana sofre menos influência de outliers. \\
84 & A & Classe dominante pode inflar acurácia; examine precision/recall e matriz de confusão. \\
85 & A & Precision = verdadeiros positivos entre todos os positivos previstos. \\
86 & A & Leakage usa informação que não estaria disponível no momento real da previsão. \\
87 & A & Correlação não demonstra causalidade. \\
88 & A & Em série temporal, mistura aleatória futuro-passado pode vazar informação. \\
89 & A & RAG recupera contexto externo e o fornece ao modelo gerador. \\
90 & A & Fine-tuning é apropriado para comportamento/padrões; RAG é melhor para conhecimento factual atualizável e rastreável. \\
91 & A & Conteúdo recuperado pode carregar instruções maliciosas e contaminar a geração. \\
92 & A & Concept drift altera a relação entre características e alvo. \\
93 & A & Ações com impacto financeiro exigem limites, autorização, logs e aprovação/segregação adequados. \\
94 & A & Explicabilidade apoia validação, entendimento e accountability, sem garantir perfeição. \\
95 & A & Suficiência = quantidade; adequação = relevância e confiabilidade/qualidade. \\
96 & A & Materialidade pode ser financeira e qualitativa, afetando escopo e julgamento. \\
97 & A & Essa cadeia estrutura um achado auditável. \\
98 & A & Amostra depende do objetivo e da população; conveniência não é regra geral. \\
99 & A & Independência e objetividade protegem julgamento profissional contra influência indevida. \\
100 & A & Monitoramento verifica implementação e efeito das providências após a auditoria. \\
\bottomrule
\end{longtable}

\subsection*{Folha de resultado}

\begin{itemize}
\item Gerais: \underline{\hspace{2cm}}/40.
\item Específicas: \underline{\hspace{2cm}}/60.
\item Nota ponderada: $G+2E=$ \underline{\hspace{2.5cm}}/160.
\item Disciplinas com desempenho abaixo de 70\%: \underline{\hspace{8cm}}.
\item Erros por interpretação: \underline{\hspace{1.5cm}}; teoria: \underline{\hspace{1.5cm}}; cálculo: \underline{\hspace{1.5cm}}; distração: \underline{\hspace{1.5cm}}.
\end{itemize}
